\chapter{Նպատակներ և սպասվող արդյունքներ}
\label{ch:goals}

\section{Նախագծի նպատակները}
\begin{enumerate}
  \item FPGA-ի ներսում իրականացնել կոնֆիգուրելի սենսոր-էմուլյացիա (RTL ազդանշան-գեներատորներ) առնվազն 3 սենսորի տիպի համար։
  \item Իրականացնել իրական ժամանակի նախնական մշակում (առնվազն 3 processing IP) FPGA-ի վրա։
  \item Ապահովել PL$\rightarrow$PS տվյալների փոխանցում և PS$\rightarrow$սերվեր հաղորդակցություն։
  \item Մշակել սերվեր, տվյալների բազա և Web հարթակ՝ պահպանման, վիզուալիզացիայի և վերլուծության համար։
  \item Ապահովել FPGA-ի \gls{runtimeconf} Web հարթակից՝ AXI4-Lite ռեգիստրների միջոցով։
  \item Ցուցադրել համակարգի ամբողջական աշխատանքը՝ փորձի (\gls{experiment}) կազմակերպման, տվյալների հավաքագրման և վերլուծության միջոցով։
\end{enumerate}

\section{Ուսումնական արդյունքներ}
Նախագծի ավարտին ուսանողները կկարողանան.
\begin{itemize}
  \item \textbf{Ուսանող 1.} նախագծել և վերիֆիկացնել RTL մոդուլներ testbench-երով, ինտեգրել դրանք AXI-ով Zynq PS-ի հետ, վերլուծել timing-ը և ռեսուրսների օգտագործումը։
  \item \textbf{Ուսանող 2.} նախագծել տվյալների մոդել և API, կառուցել Web dashboard իրական ժամանակի տվյալներով, կիրառել վիճակագրական վերլուծություն և անոմալիաների հայտնաբերում։
  \item \textbf{Երկուսն էլ.} համաձայնեցնել միջերեսներ, աշխատել git-ով, ինտեգրել և ցուցադրել ամբողջական համակարգ։
\end{itemize}
